% This document was generated by an LLM.

\documentclass{essay}

\title{Continuity, Piecewise Continuity, and Discontinuities\\
An Introduction from the Perspective of Real Analysis}
\author{}
\date{}

\begin{document}

\maketitle
\tableofcontents

%%%%%%%%%%%%%%%%%%%%%%%%%%%%%%%%%%%%%%%%%%%%%%%%%%%%%%

\section{Introduction}

One of the central goals of real analysis is to understand how functions
behave. Some functions vary smoothly, while others exhibit abrupt changes,
oscillations, or singularities. The concept of continuity provides a precise
mathematical language for distinguishing these behaviors.

Although continuity is introduced early in calculus, its true significance
becomes apparent only in analysis. Many of the fundamental theorems of calculus
are, at their core, statements about continuous functions. Conversely, studying
the different ways in which continuity can fail reveals why certain hypotheses
are necessary.

The purpose of this chapter is to develop both an intuitive and formal
understanding of

\begin{itemize}
\item continuous functions,
\item discontinuities,
\item piecewise continuous functions,
\item and their role in integration and differential equations.
\end{itemize}

Throughout, we work with functions

\[
f:D\rightarrow\mathbb{R},
\]

where the domain \(D\subseteq\mathbb{R}\).

%%%%%%%%%%%%%%%%%%%%%%%%%%%%%%%%%%%%%%%%%%%%%%%%%%%%%%

\section{Limits: The Foundation of Continuity}

Continuity is defined in terms of limits. Consequently, the notion of a limit
must come first.

\begin{definition}[Limit]

Let \(a\) be an accumulation point of the domain of \(f\).

We say

\[
\lim_{x\to a}f(x)=L
\]

if for every \(\varepsilon>0\) there exists a \(\delta>0\) such that

\[
0<|x-a|<\delta
\quad\Longrightarrow\quad
|f(x)-L|<\varepsilon.
\]

\end{definition}

Notice that the value \(f(a)\) does not appear in the definition.

A limit depends only on the behavior of the function arbitrarily close to the
point.

This distinction is one of the most important ideas in elementary analysis.

%%%%%%%%%%%%%%%%%%%%%%%%%%%%%%%%%%%%%%%%%%%%%%%%%%%%%%

\section{Continuous Functions}

The notion of continuity is obtained by requiring that the function value agree
with its limiting value.

\begin{definition}[Continuity at a Point]

A function \(f\) is continuous at \(a\) if

\[
\lim_{x\to a}f(x)=f(a).
\]

\end{definition}

Intuitively, nearby inputs produce nearby outputs.

The function behaves predictably under arbitrarily small perturbations of the
input.

\begin{definition}[Continuous on an Interval]

A function is continuous on an interval if it is continuous at every point of
that interval.

\end{definition}

%%%%%%%%%%%%%%%%%%%%%%%%%%%%%%%%%%%%%%%%%%%%%%%%%%%%%%

\section{Equivalent Characterizations}

Real analysis offers several equivalent ways to express continuity.

\begin{theorem}

For a function \(f:\mathbb{R}\rightarrow\mathbb{R}\), the following are
equivalent.

\begin{enumerate}
\item \(f\) is continuous at \(a\).

\item Whenever \(x_n\to a\),

\[
f(x_n)\to f(a).
\]

\item The inverse image of every open set is open (topological definition).

\end{enumerate}

\end{theorem}

The second characterization is often easier to use in proofs, while the third
extends naturally to general metric and topological spaces.

%%%%%%%%%%%%%%%%%%%%%%%%%%%%%%%%%%%%%%%%%%%%%%%%%%%%%%

\section{How Continuity Can Fail}

Discontinuities arise whenever

\[
\lim_{x\to a}f(x)\neq f(a)
\]

or the limit does not exist.

However, the failure may occur in several fundamentally different ways.

%%%%%%%%%%%%%%%%%%%%%%%%%%%%%%%%%%%%%%%%%%%%%%%%%%%%%%

\subsection{Removable Discontinuities}

\begin{example}

Consider

\[
f(x)=\frac{x^2-1}{x-1},
\qquad x\neq1.
\]

Since

\[
\frac{x^2-1}{x-1}=x+1,
\]

we obtain

\[
\lim_{x\to1}f(x)=2.
\]

The limit exists, but the function is not defined at \(x=1\).

By defining

\[
f(1)=2,
\]

the discontinuity disappears.

\end{example}

A removable discontinuity represents missing information rather than genuinely
irregular behavior.

%%%%%%%%%%%%%%%%%%%%%%%%%%%%%%%%%%%%%%%%%%%%%%%%%%%%%%

\subsection{Jump Discontinuities}

A jump discontinuity occurs when the one-sided limits both exist but differ.

\begin{example}

\[
f(x)=
\begin{cases}
0,&x<0,\\
1,&x\ge0.
\end{cases}
\]

Then

\[
\lim_{x\to0^-}f(x)=0,
\qquad
\lim_{x\to0^+}f(x)=1.
\]

Since these limits are unequal, the ordinary limit does not exist.

\end{example}

%%%%%%%%%%%%%%%%%%%%%%%%%%%%%%%%%%%%%%%%%%%%%%%%%%%%%%

\subsection{Infinite Discontinuities}

Infinite discontinuities occur when the function grows without bound.

The simplest example is

\[
f(x)=\frac1x.
\]

Near the origin,

\[
|f(x)|\rightarrow\infty.
\]

%%%%%%%%%%%%%%%%%%%%%%%%%%%%%%%%%%%%%%%%%%%%%%%%%%%%%%

\subsection{Oscillatory Discontinuities}

Perhaps the most important counterexample in introductory analysis is

\[
f(x)=
\sin\left(\frac1x\right).
\]

As \(x\rightarrow0\), the oscillations become infinitely frequent.

No limiting value exists.

Unlike a jump discontinuity, there is no preferred left or right value.

%%%%%%%%%%%%%%%%%%%%%%%%%%%%%%%%%%%%%%%%%%%%%%%%%%%%%%

\section{One-Sided Limits}

Many practical functions possess jumps.

To describe them, we introduce one-sided limits.

\begin{definition}

The left-hand limit is

\[
\lim_{x\to a^-}f(x),
\]

while the right-hand limit is

\[
\lim_{x\to a^+}f(x).
\]

\end{definition}

\begin{theorem}

The ordinary limit

\[
\lim_{x\to a}f(x)
\]

exists if and only if

\[
\lim_{x\to a^-}f(x)
=
\lim_{x\to a^+}f(x).
\]

\end{theorem}

%%%%%%%%%%%%%%%%%%%%%%%%%%%%%%%%%%%%%%%%%%%%%%%%%%%%%%

\section{Piecewise Continuous Functions}

Many functions encountered in applications are continuous except at isolated
jump discontinuities.

\begin{definition}

A function is piecewise continuous on a closed interval \([a,b]\) if

\begin{enumerate}
\item it has only finitely many discontinuities;

\item on every open subinterval between consecutive discontinuities it is
continuous;

\item every discontinuity has finite left and right limits.

\end{enumerate}

\end{definition}

For functions defined on

\[
[0,\infty),
\]

the definition requires every finite interval

\[
[0,b]
\]

to satisfy the above conditions.

Consequently, infinitely many discontinuities are allowed over the entire
half-line provided they do not accumulate inside any finite interval.

%%%%%%%%%%%%%%%%%%%%%%%%%%%%%%%%%%%%%%%%%%%%%%%%%%%%%%

\section{Examples}

\begin{example}

Every polynomial is continuous.

\end{example}

\begin{example}

Every rational function is continuous wherever its denominator is nonzero.

\end{example}

\begin{example}

The floor function

\[
\lfloor x\rfloor
\]

is piecewise continuous.

It has jump discontinuities at every integer.

\end{example}

\begin{example}

The sign function

\[
\operatorname{sgn}(x)
\]

is piecewise continuous.

\end{example}

\begin{example}

The function

\[
\sin\left(\frac1x\right)
\]

is not piecewise continuous on any interval containing the origin.

\end{example}

%%%%%%%%%%%%%%%%%%%%%%%%%%%%%%%%%%%%%%%%%%%%%%%%%%%%%%

\section{Important Theorems About Continuous Functions}

The following theorems form the backbone of elementary analysis.

Proofs are omitted.

\begin{theorem}[Algebra of Continuous Functions]

Sums, differences, products, and scalar multiples of continuous functions are
continuous.

Whenever defined, quotients of continuous functions are continuous.

\end{theorem}

\begin{theorem}[Composition]

The composition of continuous functions is continuous.

\end{theorem}

\begin{theorem}[Extreme Value Theorem]

Every continuous function on a compact interval

\[
[a,b]
\]

attains both a maximum and a minimum.

\end{theorem}

\begin{theorem}[Intermediate Value Theorem]

If

\[
f(a)<c<f(b),
\]

then there exists

\[
x\in(a,b)
\]

such that

\[
f(x)=c.
\]

\end{theorem}

\begin{theorem}[Uniform Continuity]

Every continuous function on a compact interval is uniformly continuous.

\end{theorem}

%%%%%%%%%%%%%%%%%%%%%%%%%%%%%%%%%%%%%%%%%%%%%%%%%%%%%%

\section{Continuity and Integration}

One reason continuity is so important is its relationship with integration.

\begin{theorem}

Every continuous function on a closed bounded interval is Riemann integrable.

\end{theorem}

The converse is false.

Many discontinuous functions are nevertheless Riemann integrable.

For example, every piecewise continuous function on a closed interval is
Riemann integrable.

Thus jump discontinuities do not destroy integrability.

%%%%%%%%%%%%%%%%%%%%%%%%%%%%%%%%%%%%%%%%%%%%%%%%%%%%%%

\section{Continuity and Differentiability}

Differentiability is a stronger condition than continuity.

\begin{theorem}

If \(f\) is differentiable at \(a\), then \(f\) is continuous at \(a\).

\end{theorem}

The converse does not hold.

The classical example is

\[
f(x)=|x|,
\]

which is continuous everywhere but not differentiable at the origin.

%%%%%%%%%%%%%%%%%%%%%%%%%%%%%%%%%%%%%%%%%%%%%%%%%%%%%%

\section{A Hierarchy of Regularity}

Functions encountered in analysis can be arranged according to increasing
regularity.

\[
\begin{aligned}
\text{arbitrary functions}
&\\
\supset
\text{Riemann integrable}
&\\
\supset
\text{piecewise continuous}
&\\
\supset
\text{continuous}
&\\
\supset
C^1
&\\
\supset
C^2
&\\
\supset
\cdots
&\\
\supset
C^\infty
&\\
\supset
\text{analytic}.
\end{aligned}
\]

Each level imposes additional structure.

As assumptions become stronger, more powerful theorems become available.

%%%%%%%%%%%%%%%%%%%%%%%%%%%%%%%%%%%%%%%%%%%%%%%%%%%%%%

\section{Why Piecewise Continuity Appears So Often}

Many physical systems exhibit abrupt but isolated changes.

Examples include

\begin{itemize}
\item electrical switches,
\item relay circuits,
\item digital signals,
\item square waves,
\item pulse trains,
\item piecewise constant forcing terms in differential equations.
\end{itemize}

These systems are not continuous, yet they remain sufficiently regular for the
machinery of calculus to apply.

Consequently, piecewise continuity is a standard hypothesis in the theory of
Laplace transforms, Fourier series, and differential equations.

%%%%%%%%%%%%%%%%%%%%%%%%%%%%%%%%%%%%%%%%%%%%%%%%%%%%%%

\section{Summary}

The concept of continuity formalizes the intuitive idea that small changes in
the input produce small changes in the output.

Discontinuities arise in several qualitatively different ways, including
removable, jump, infinite, and oscillatory discontinuities.

Among discontinuous functions, piecewise continuous functions occupy an
especially important position. They preserve many of the useful analytical
properties of continuous functions while allowing isolated jumps that naturally
occur in applications.

Understanding continuity is therefore not merely an exercise in manipulating
limits. It is the first step toward understanding why the major theorems of
analysis are true and precisely which hypotheses they require.

\end{document}
